\section{Conclusion}

We document a fundamental structural shift in options markets: the emergence of persistent dealer gamma regimes coinciding with 0DTE proliferation. Five-phase validation across 2,221 total windows (1,412 real market windows plus 809 synthetic controls) spanning six years (2020-2025) reveals that sustained dealer positioning---characterized by persistence ($\geq$70\% same-sign dominance), magnitude ($\geq$\$5B average), and stability ($\leq$5 sign flips)---has transitioned from an exceptional state (12\% detection in 2020) to a permanent characteristic of post-0DTE markets (100\% in 2024).

\subsection{Key Findings}

\textbf{Framework selectivity validated}: The 69.1 percentage point discrimination between 2024 (81.2\% detection) and 2020 (12.1\% detection) provides evidence that our methodology can detect persistent structural regimes, not universal patterns. This 5.7x separation (p $<$ 0.0001, $\phi$ = 0.672) suggests the framework effectively distinguishes sustained dealer constraints from fragmented markets.

\textbf{Negative controls passed}: Phase 2 achieved 0\% false positives on transitional (7-10 flips) and low-magnitude ($<$\$5B) synthetic windows, providing evidence that the framework enforces regime criteria. The 12.1\% false positive rate on 2020 shuffled data further validates that detection depends on temporal structure, not statistical artifacts.

\textbf{Gradual 0DTE adoption pattern identified}: Phase 5 multi-year validation (1,412 windows, 2020-2025) revealed that the structural shift occurred through gradual adoption culminating in a 2023→2024 transition, not a sharp 2020→2021 change. Detection rates closely tracked market evolution: 2020 (12.2\%), 2021 (3.7\%), 2022 (32.4\%), 2023 (20.2\%), 2024 (100\%), 2025 (100\%). This gradual pattern is consistent with the mechanistic interpretation—detection tracks 0DTE market penetration rather than exhibiting binary behavior.

\textbf{Magnitude-driven regime shift}: GEX magnitude grew 360\% from \$5B (2021) to \$23B (2024), far exceeding 20-25\% cumulative inflation. The fixed \$5B threshold effectively discriminated borderline markets (2021: 3.7\% detection) from structural regimes (2024: 100\% detection), supporting threshold design. The 4.95× detection increase (20.2\% → 100\%) temporally coincides with 0DTE trading explosion in late 2023.

\textbf{LLM reasoning quality}: Manual review of 50 randomly sampled responses revealed 98\% mechanical accuracy, 100\% criterion citation, and 88\% step-by-step verification across both detection conditions, suggesting structural reasoning rather than pattern matching.

\textbf{Robustness against inflation}: Price normalization experiments confirm that the 2020→2024 detection increase is not an artifact of nominal asset growth. Applying the SPY price growth factor (2.3$\times$) to 2020 GEX magnitudes accounts for only 31\% of the detection gap; 69\% persists after normalization due to structural differences in persistence (96\% vs.\ 69\%) and stability (99\% vs.\ 54\%). The framework discriminates 0DTE-driven regime persistence from pre-2021 dynamics even when controlling for capitalization changes.

\textbf{Calculation independence}: Formula Agreement Testing confirms that regime detection is robust to mathematical formulation. The LLM achieved 100\% agreement (61/61 windows) between absolute dollar-scaled GEX and normalized GEX ratios, demonstrating the plausibility that the model identifies regimes based on latent structural positioning (persistence and stability) rather than relying solely on nominal magnitude cues. This complements the inflation defense: 2020 fails even with inflated magnitude (lacks structure), while 2024 succeeds even without magnitude (has structure).

\subsection{Contributions to Market Microstructure Analysis}

This work advances understanding of post-0DTE market structure in four ways:

\begin{enumerate}
\item \textbf{Temporal obfuscation methodology}: First application of obfuscation testing to multi-day regimes (30-day windows vs single-day snapshots), establishing that LLMs can reason about sustained constraints across pre-0DTE (2020) and post-0DTE (2024) market eras
\item \textbf{Selectivity validation framework}: Comprehensive negative controls demonstrating $<$10\% false positive rate, addressing a key limitation of prior LLM finance validation that typically tests only positive cases
\item \textbf{Cross-era regime comparison}: Phase 4 isolated a large 2020→2024 structural difference (12.1\% → 81.2\% detection, 69.1pp, $\phi$ = 0.672), quantifying the market structure change between pre-0DTE and post-0DTE eras
\item \textbf{Scalable validation methodology}: OpenAI Batch API enabled comprehensive five-phase validation across 2,221 windows, making rigorous multi-year testing computationally feasible for academic research
% Cost details (\$11.67, 50\% savings vs synchronous) omitted from paper
\end{enumerate}

\subsection{Implications}

Our findings have three implications for market microstructure research:

\textbf{Temporal reasoning capabilities}: LLMs can identify persistent structural constraints across 30-day windows when given appropriate prompts and criteria. Manual review of 50 randomly sampled windows revealed 98\% mechanical accuracy, suggesting robust computational reasoning about dealer hedging mechanics across all 2,221 validation windows.

\textbf{Selectivity is essential}: Our finding that 5-day windows achieved 98-100\% detection (too universal) while 30-day windows with strict criteria achieved 12-81\% detection (selective) underscores that validation frameworks must test discrimination, not just positive case accuracy. Future work should include cross-regime comparison and negative controls.

\textbf{Structural vs statistical patterns}: High LLM mechanical accuracy (98\%) combined with 0\% false positives on synthetic tests demonstrates the model reasons about dealer hedging constraints, not spurious statistical patterns. This structural focus sidesteps p-hacking risks inherent in profitability-based validation.

\subsection{Future Directions}

Four research directions emerge from this work:

\textbf{Historical extension}: While Phase 5 covered the 2020-2025 0DTE transition, longer-term historical analysis (2005-2020) could reveal whether similar regime transitions occurred during earlier options market innovations such as 2008-2010 (financial crisis), 2016-2018 (VIX ETF proliferation), and 2005-2007 (pre-crisis baseline).

\textbf{Cross-asset generalization}: Extending to individual equities (AAPL, MSFT, NVDA, TSLA) will test whether regime detection generalizes beyond S\&P 500 index options and identify asset-specific vs market-wide regime characteristics.

\textbf{Volume-based GEX measurement}: Incorporating intraday options flow data (when available) may refine regime detection by capturing 0DTE hedging dynamics missed by end-of-day open interest measurements.

\textbf{Tiered regime classification}: Future work should develop tiered classifications (moderate/strong/extreme regimes) to discriminate intensity within persistent regime states. LLM confidence scores already correlate significantly with regime quality ($r = +0.501$ for persistence, $r = -0.549$ for stability), suggesting a natural basis for intensity grading beyond binary detection.

\subsection{Concluding Remarks}

This paper documents the emergence of persistent dealer gamma regimes as a distinguishing characteristic of post-0DTE options markets. The 69.1 percentage point detection difference between 2024 (81.2\%) and 2020 (12.1\%) quantifies a fundamental market structure change---what we term the ``scar tissue'' of 0DTE trading, where concentrated intraday hedging flows leave residual inventory that constrains overnight positioning.

The ``great divergence'' in SPY returns provides supporting evidence: overnight returns (+21.2\%) have significantly outperformed intraday returns (+0.6\%) since late 2023, consistent with a displacement of risk premium to the overnight session. Our structural identification framework, validated through comprehensive negative controls and cross-era comparison, offers a methodology for detecting sustained dealer constraints and tracking market structure evolution as 0DTE trading continues to reshape options market microstructure.
